\chapter{Proemio}

\textit{Ai discepoli, ai maestri. Ai curiosi e ai volenterosi.}

Noi bardi, guerrieri, maghi, druidi, chierici e avventurieri, ci riuniamo alla fine delle nostre vite, per ultimare il fluire del nostro scopo, nel divulgare la dottrina che ci ha guidato. Raccoglieremo le nostre esperienze e le nostre conoscenze per tratteggiare gli insegnamenti dello scorrere del fiume, il cui scopo è vivere una vita colma di un significato autentico, creativo e impegnativo, secondo il proprio consapevole sentire.
Noi cancelleremo i nostri nomi e le nostre immagini, per quanto ci è possibile, dalla storia, dal momento che non desideriamo essere né ricordati né venerati. Il nostro desiderio consapevole è offrire ai futuri fluminauti gli strumenti per modellare il corso del proprio fiume interiore, e tale trasformazione, continua e inarrestabile, non può sottendere al volere di alcun dio, al capriccio di alcun re, alla regola di alcun legislatore, al dogma di alcuna setta. La sola autorità dello scorrere è quella di chi abbraccia la dottrina volontariamente e con sincera curiosità, spinto dal proprio sentire. Siate forti e pazienti, perché l’impegno sarà molto, costante e continuo. E se a un punto vorrete affrancarvene non siate duri con voi stessi, giacché noi non vogliamo impedire il vario e numeroso scorrere dei fiumi di questo vasto mondo.

Lo scopo di un individuo è sempre morale. Dal più piccolo gesto, come bere da una ciotola, alle più grandi imprese, come fondare una nazione, il vero scopo è nell'attuazione stessa e non nella risoluzione finale. La causa da cui scaturisconon le azioni della nostra vita hanno sempre a che fare con una spinta interiore, una forza che non possiamo controllare ma che deve trovare sfogo. Questa forza vitale consente all'individuo di esistere ma, come un fiume, l'individuo che la naviga non può fermarla, giacché fermarla significa morire. Noi chiamiamo questa forza lo "scorrere" o il "fluire del proprio fiume.

Se un undividuo si lascia trascinare dal proprio fiume, egli agirà secondo morale, ma tale morale sarà vaga e caotica, ed egli non capirà la ragione della proproa scelta, non capirà i veri motivi del suo sentirsi amareggiato, deluso, rabbioso. I momenti di felicità, euforia, serenità, saranno vissuti con trasporto, ma la loro utilità si esaurirà al loro termine, senza fecondità per il futuro. Costui sarà libero, ma la sua libertà sarà quella di una foglia nel vento.

Ma se costui capirà le proprie emozioni nel loro intero susseguirsi, e le osserverà comporre un quadro, potrà scoprire se tale disegno è coerente e nel caso no lo sia, potrà intravedere il quelle incoerenze uno scopo, il proprio fluire, teso in una direzione che sfugge alle apparenze. Egli potrà riflettere se tale scorrere è desiderabile e, attingendo alle proprie conoscenze, potrà vedere chiaramente uno scopo. Allora il suo fluire campbierà in accordo con questa prassi che noi chiamiamo "psicotrofia". Egli non sarà totalmente in balia del suo fluire e potrà vivere le proprie emozioni in modo pieno, privo di sensi di colpa o smarrimento che ottunda il giudizio, e senza confondere queste con la propria identità, giacchè le emozioni sono i mediatori fra la realtà e il sé, e l'identità non è un incisione immutabile ma è un esperienza continua del proprio divenire.
